% Adaptado por Andre
% ============================================================
% PLANTILLA JAES -- ESTILO DE CITACIÓN: APA 7 (biblatex-apa, oficial)
% ============================================================
% Citas como (Autor, Año) / Autor (Año), bibliografía alfabetizada,
% generadas con la implementación OFICIAL de APA 7 para LaTeX:
% el paquete biblatex-apa (https://github.com/plk/biblatex-apa).
%
% A diferencia de la aproximación con natbib+apalike.bst, esta
% plantilla usa el motor BIBLATEX (no BibTeX clásico), que requiere
% el backend BIBER en vez de bibtex. Esto es un cambio de flujo de
% compilación, no solo de paquete -- ver notas de compilación abajo.
%
% ---- Entorno de referencia de esta plantilla ----
%   SO:           Windows 11
%   Editor:       VS Code + extensión LaTeX Workshop
%   Motor:        latexmk (vía Strawberry Perl)
%   Distribución: MiKTeX, con auto-instalación de paquetes activa
%
% ---- Estado de esta plantilla ----
% Verificada por compilación completa (Jun-2026): 0 errores, 0 warnings,
% biber sin errores ni warnings, citas \textcite/\parencite funcionando
% con formato APA 7 real (incluye "&" en cita parentética con dos
% autores, y sentence case automático en los títulos de la bibliografía).
%
% Se encontró y corrigió un conflicto real entre jaes.cls y biblatex:
% jaes.cls define \bibparsep y \biblabelsep con \newdimen (para su
% entorno thebibliography clásico); biblatex declara esos MISMOS
% nombres con \newlength sin comprobar si ya existen, lo que rompía la
% compilación con "Command \bibparsep already defined". La solución
% (liberar ambos nombres antes de cargar biblatex) ya está aplicada
% más abajo en este archivo -- no necesitas hacer nada al respecto.
%
% Archivos que deben estar en esta misma carpeta:
%   - plantilla_apa7.tex  (este archivo)
%   - jaes.cls            (clase del documento, define el formato JAES)
%   - jaes.bib            (tu base de datos de referencias, formato
%                          biblatex -- ver comentarios dentro del .bib)
%   (NO necesitas ningún .bst aquí; biblatex-apa no usa archivos .bst,
%    todo el formato lo controla el paquete biblatex-apa vía biber)
%
% Compilación con LaTeX Workshop + latexmk:
%   1. Abre esta carpeta en VS Code.
%   2. La PRIMERA vez que compiles, MiKTeX detectará que faltan los
%      paquetes biblatex, biblatex-apa, csquotes y biber, y te
%      preguntará si quieres instalarlos. Acepta -- se instalan solos,
%      sin que tengas que buscar nada manualmente.
%   3. LaTeX Workshop necesita saber que este documento usa biber, no
%      bibtex. Por defecto, la receta estándar de LaTeX Workshop
%      ("latexmk") ya detecta esto automáticamente leyendo el .log
%      (busca la línea "Backend: biber" generada por biblatex). Si tu
%      compilación se queda esperando o no actualiza las citas, revisa
%      en VS Code: Configuración → busca "latex-workshop.latex.recipes"
%      y confirma que la receta por defecto sea la que incluye
%      "latexmk" (no "pdflatex" suelto), o agrega manualmente esta
%      receta en tu settings.json del proyecto:
%        "latex-workshop.latex.recipes": [
%          { "name": "latexmk", "tools": ["latexmk"] }
%        ],
%        "latex-workshop.latex.tools": [
%          { "name": "latexmk", "command": "latexmk",
%            "args": ["-synctex=1","-interaction=nonstopmode",
%                     "-file-line-error","-pdf","-outdir=%OUTDIR%","%DOC%"] }
%        ]
%      (latexmk internamente decide solo si debe llamar a bibtex o a
%      biber, según lo que detecte en el .aux -- no necesitas elegir).
%   4. Si compilas por terminal en vez de con el botón de VS Code:
%        latexmk -pdf plantilla_apa7.tex
%      (un solo comando; latexmk corre pdflatex/biber las veces que
%      haga falta automáticamente, a diferencia del flujo manual de
%      bibtex de 4 pasos que se usa con jaes.bst/apalike.bst).
%
% ---- Solución de problemas ----
%   - Error "! LaTeX Error: Command \bibparsep already defined" (o lo
%     mismo con \biblabelsep) → ya está resuelto en esta plantilla (ver
%     nota arriba). Si lo ves de todas formas, revisa que no hayas
%     borrado por accidente las dos líneas \let...\undefined que
%     aparecen justo antes de \usepackage{biblatex} más abajo.
%   - Error "Package biblatex Error: Missing \bibliography...command"
%     → falta el \addbibresource{jaes.bib} en el preámbulo (revisa que
%       no lo hayas borrado por accidente).
%   - Citas como "??" o el aviso "Please (re)run Biber on the file" tras
%     compilar una sola vez → ESTO ES NORMAL en la primera compilación
%     de un documento nuevo: biblatex necesita generar primero un
%     archivo .bcf (lo hace en la 1.ª pasada de pdflatex), luego biber
%     lee ese .bcf y genera el .bbl, y recién entonces pdflatex puede
%     mostrar las citas correctas. En VS Code esto se resuelve solo
%     pulsando el botón ▶ de compilar una segunda vez (LaTeX Workshop
%     vuelve a correr pdflatex → biber → pdflatex → pdflatex). Si
%     después de compilar 2-3 veces el problema persiste, revisa que
%     LaTeX Workshop esté usando la receta "latexmk" (ver punto 3 de
%     compilación arriba), no una receta manual de solo "pdflatex".
%   - "and is already a type wizard-defined function name" o errores
%     parecidos NO deberían aparecer aquí (esos eran específicos de
%     jaes.bst con BibTeX clásico); si aparecen, probablemente VS Code
%     está intentando correr bibtex en vez de biber sobre este
%     proyecto -- revisa el punto 3 de compilación arriba.
%   - Si el formato de la bibliografía sale en inglés en vez de
%     español (ej. "and" en vez de "y" entre dos autores en cita
%     NARRATIVA con \textcite -- la cita PARENTÉTICA con \parencite ya
%     usa "&" correctamente sin necesitar nada adicional), agrega
%     \usepackage[spanish]{babel} en el preámbulo: biblatex-apa SÍ trae
%     localización al español lista (american-apa.lbx existe junto a
%     otros idiomas), pero solo se activa si babel está cargado con
%     ese idioma.
% ============================================================

\documentclass[fleqn, final]{jaes}

% ---- Metadatos del pie de página ----
\jyear{2026}                          % Año que aparece en el pie de página
\jmonth{Junio}                        % Mes que aparece en el pie de página
\jname{Nombre del Curso -- Código}     % Texto del pie (por defecto: "Proyecto de Investigación")

% ---- Paquetes adicionales ----
\usepackage[utf8]{inputenc}

% csquotes es requisito de biblatex-apa (>= 4.3). Se autoinstala con
% MiKTeX la primera vez que compiles.
\usepackage{csquotes}

% FIX VERIFICADO (Jun-2026): jaes.cls ya define \bibparsep y
% \biblabelsep con \newdimen (para su propio entorno thebibliography
% clásico, el que usan las plantillas IEEE/apalike). biblatex también
% declara esos mismos nombres con \newlength sin comprobar si ya
% existen, lo que produce "! LaTeX Error: Command \bibparsep already
% defined" al cargar el paquete. Se liberan los nombres justo antes
% de \usepackage{biblatex} para evitar el choque; esto no afecta el
% resto de jaes.cls porque esas dos longitudes solo se usan dentro del
% entorno thebibliography clásico, que esta plantilla no utiliza
% (usamos \printbibliography de biblatex en su lugar).
\let\bibparsep\undefined
\let\biblabelsep\undefined

% biblatex con estilo APA 7 oficial. backend=biber es obligatorio
% (biblatex-apa no funciona con el backend bibtex clásico).
\usepackage[backend=biber, style=apa]{biblatex}
\addbibresource{jaes.bib}   % tu base de datos de referencias (con extensión .bib aquí)

\usepackage{amsmath}\setlength{\mathindent}{10pt}
\usepackage{bm}

% cleveref: permite \cref{} / \Cref{} para referencias automáticas
% ("Sec. 3", "Tabla 2", etc.) sin escribir el número a mano.
\usepackage[capitalise]{cleveref}
  \crefformat{section}{\textsc{Sec.~#2#1#3}}
  \Crefformat{section}{\textsc{Sec.~#2#1#3}}
  \crefformat{figure}{{Fig.~#2#1#3}}
  \Crefformat{figure}{{Fig.~#2#1#3}}
  \crefformat{table}{{Tabla~#2#1#3}}
  \Crefformat{table}{{Tabla~#2#1#3}}
  \crefname{section}{Sec.}{Secs.}
  \Crefname{section}{Sec.}{Secs.}
  \crefname{table}{Tabla}{Tablas}
  \Crefname{table}{Tabla}{Tablas}
  \crefname{figure}{Fig.}{Figs.}
  \Crefname{figure}{Fig.}{Figs.}
% Fix: cleveref dice "and" en vez de "y" al combinar 2+ referencias
% (\cref{sec:a,sec:b}) si el documento no usa el paquete babel.
% Sin este bloque, saldría "(Sec. 2 and 5)" en inglés.
\providecommand{\creflastconjunction}{}
\renewcommand{\creflastconjunction}{ y }
\providecommand{\crefpairconjunction}{}
\renewcommand{\crefpairconjunction}{ y }
\providecommand{\crefmiddleconjunction}{}
\renewcommand{\crefmiddleconjunction}{, }

% Paquetes para tablas anchas / con texto que necesita ajustarse de línea
% (muy recomendable si vas a tener tablas con párrafos largos en celdas).
\usepackage{array}
\usepackage{tabularx}
\usepackage{longtable}
\usepackage{multirow}
\usepackage{ragged2e}
\newcolumntype{Y}{>{\RaggedRight\arraybackslash}X}  % columna de ancho fijo con ajuste de línea

% Localización al español (jaes.cls viene en inglés por defecto):
\renewcommand{\tablename}{Tabla}      % por defecto decía "Table"
% \bibname (NO \refname) es el comando que jaes.cls usa internamente
% para titular secciones de bibliografía con el entorno clásico
% thebibliography. biblatex usa su PROPIO título ("Referencias" o
% "References" según idioma) vía \printbibliography, que no depende
% de \bibname -- esta línea se deja solo por consistencia/seguridad,
% no tiene efecto visible en esta plantilla.
\renewcommand{\bibname}{Referencias}

% Anula la biografía automática de autores (\biography), que tiene un
% bug conocido en esta clase: no imprime el texto biográfico, solo
% fotos y nombres. Ver la sección "Sobre el Equipo de Trabajo" más abajo
% para la alternativa manual que SÍ funciona.
\def\printbiography{}

% ---- Si tu entrega NO debe decir "BORRADOR" en cada página ----
% jaes.cls imprime la palabra BORRADOR en el encabezado de todas las
% páginas de forma fija (no es una opción togglable desde el .tex).
% Para quitarlo hay que editar tu copia de jaes.cls, buscando estas
% dos líneas y dejando vacío el contenido de \rhfont:
%
%   \def\@evenhead{\hbox to \textwidth{\hfill\rhfont BORRADOR}}%
%   \def\@oddhead{\hbox to \textwidth{\rhfont BORRADOR\hfill}}%
%
% (Si tu archivo aún es un borrador en revisión, puedes dejarlo así).

\begin{document}

% Encabezado de cada página: {Lista de autores}{Título corto del trabajo}
\markboth{APELLIDO1, APELLIDO2, APELLIDO3}{Título Corto del Trabajo}

% ---------------- Título ----------------
\title{Título Completo del Trabajo}

% ---------------- Autores ----------------
% Un \author{...} por persona, separados por comas, todos dentro de
% un único \authorgroup{...}. El \textsuperscript{1} enlaza cada autor
% con su afiliación correspondiente en \affil{}.
% IMPORTANTE: \email{} y \affil{} deben ir DENTRO del bloque
% \authorgroup{...} (antes de la llave de cierre), no después. Si se
% colocan fuera, su formato (cursiva, tamaño de fuente) se filtra al
% resto del documento porque \authorgroup solo GUARDA su contenido
% para imprimirlo más tarde con \maketitle; lo que esté fuera de su
% llave de cierre se ejecuta de inmediato, en un punto del documento
% donde el diseño de página todavía no está listo para recibirlo.
\authorgroup{
\author{NOMBRE APELLIDO\textsuperscript{1}\orcidlink{0000-0000-0000-0000}},
\author{NOMBRE APELLIDO 2\textsuperscript{1}},
\author{NOMBRE APELLIDO 3\textsuperscript{1}}
% \email{...} es opcional: úsalo para correos, códigos de alumno, etc.
\email{(código1) \quad (código2) \quad (código3)}
\affil{\textsuperscript{1}Universidad, Facultad, Carrera, Ciudad, País}
}

% ---------------- Resumen ----------------
% El abstract va en un solo párrafo, sin saltos de línea en blanco.
% Longitud orientativa: 150-250 palabras. Debe poder leerse de forma
% independiente del resto del documento (qué se hizo, cómo, y qué se
% encontró), sin citar tablas/figuras por número.
\abstract{%
Texto del resumen aquí. Describe el problema u objetivo, la metodología
utilizada, los principales hallazgos y la conclusión o recomendación
central del trabajo, en un solo párrafo corrido.}
\maketitle

% ---------------- Palabras clave ----------------
\begin{keywords}
palabra clave 1, palabra clave 2, palabra clave 3, palabra clave 4
\end{keywords}

% ============================================================
% CUERPO DEL DOCUMENTO
% Reemplaza estos títulos de sección por los que pida tu rúbrica/curso.
% \label{sec:...} permite referenciar la sección con \cref{sec:...}
% desde cualquier otra parte del documento.
% ============================================================

\section{Introducción / Diagnóstico}\label{sec:introduccion}

Texto de la sección. Usa \verb|\subsection{...}| para subdividir si
la sección es larga.

\subsection{Subsección de ejemplo}

Para citar una fuente del archivo .bib en estilo APA con biblatex, usa:
\begin{itemize}
\item \verb|\textcite{clave}| → cita narrativa: Autor (Año)
\item \verb|\parencite{clave}| → cita entre paréntesis: (Autor, Año)
\item \verb|\parencite{clave1,clave2}| → varias fuentes: (Autor1, Año; Autor2, Año)
\item \verb|\textcite{clave}| seguido de página: \verb|\parencite[p.~5]{clave}| → (Autor, Año, p. 5)
\end{itemize}
% Nota: con natbib se usaría \citep/\citet; con biblatex (este caso)
% los comandos correctos son \parencite y \textcite. Son parecidos
% pero NO intercambiables entre paquetes.
Ejemplo: según \textcite{ejemplo_clave}, el hallazgo principal fue...
\parencite{ejemplo_clave}.

% ---- Ejemplo de tabla simple (una columna) ----
\begin{table}[tb]
\tabcolsep8.1pt
\centering
\caption{Título descriptivo de la tabla.}
\label{table:ejemplo}
{%
\begin{tabular}{@{}lcc@{}}\toprule
Columna 1 & Columna 2 & Columna 3\\\colrule
Fila 1 & dato & dato\\
Fila 2 & dato & dato\\\botrule
\end{tabular}}
\end{table}

Se referencia así: ver \cref{table:ejemplo}.

% ---- Ejemplo de tabla ancha (a dos columnas, usa tabularx) ----
% Usar table* (no table) para que ocupe el ancho completo de la página.
% La columna tipo Y permite que el texto largo haga salto de línea
% automático en vez de desbordarse fuera de la página.
\begin{table*}[tb]
\centering
\caption{Título descriptivo de la tabla ancha.}
\label{table:ejemplo_ancho}
{\small
\begin{tabularx}{\textwidth}{@{}lYc@{}}\toprule
Columna corta & Columna de texto largo (ajusta línea) & Columna corta\\\colrule
Dato & Texto largo que se ajustará automáticamente dentro del ancho de columna asignado, sin desbordarse. & Dato\\\botrule
\end{tabularx}}
\end{table*}

\section{Metodología}\label{sec:metodologia}

Texto de la sección.

\section{Resultados / Análisis}\label{sec:resultados}

Texto de la sección.

\section{Conclusiones}\label{sec:conclusiones}

Texto de la sección.

% ============================================================
% Bibliografía -- ESTILO APA 7 OFICIAL (biblatex-apa)
% ============================================================
% \printbibliography reemplaza por completo a \bibliography{}+
% \bibliographystyle{} del flujo BibTeX clásico. No necesita el
% nombre del archivo .bib aquí porque ya se indicó arriba con
% \addbibresource{jaes.bib} en el preámbulo.
\printbibliography[title={Referencias}]

% ============================================================
% Equipo de trabajo (alternativa manual a \biography, que no
% imprime el texto biográfico correctamente en esta clase)
% ============================================================
\clearpage
\markboth{Equipo de Trabajo}{Equipo de Trabajo}
\section*{Sobre el Equipo de Trabajo}

\begin{description}
\item[Nombre Apellido (código)] Rol o aporte específico al trabajo.
\item[Nombre Apellido 2 (código)] Rol o aporte específico al trabajo.
\item[Nombre Apellido 3 (código)] Rol o aporte específico al trabajo.
\end{description}

\noindent Texto de cierre opcional: institución, curso, docente, etc.

\end{document}
